\documentclass[10pt]{article}

% ── Packages ──────────────────────────────────────────────────────────────────
\usepackage[margin=1.1in,top=1in,bottom=1in]{geometry}
\usepackage{graphicx}
\usepackage{booktabs}
\usepackage{amsmath,amssymb}
\usepackage{hyperref}
\usepackage[numbers,sort&compress]{natbib}
\usepackage{xcolor}
\usepackage{microtype}
\usepackage{caption}
\usepackage{subcaption}
\usepackage{float}           % [H] float placement
\usepackage{tabularx}
\usepackage{multirow}
\usepackage{array}
\usepackage{parskip}         % paragraph spacing instead of indent in single-col

\floatplacement{figure}{H}   % default: place figure exactly here

\hypersetup{
  colorlinks=true,
  linkcolor=blue!70!black,
  citecolor=blue!70!black,
  urlcolor=blue!70!black
}

\graphicspath{{/data/ross/interp/paper/figures/}{/data/ross/interp/latent_analysis/}}

\captionsetup{font=small, labelfont=bf, skip=5pt}

% ── Macros ────────────────────────────────────────────────────────────────────
\newcommand{\prot}[1]{\texttt{#1}}
\newcommand{\latent}[1]{$\ell_{#1}$}
\newcommand{\Zn}{Zn$^{2+}$}
\newcommand{\Ca}{Ca$^{2+}$}
\newcommand{\Mg}{Mg$^{2+}$}
\newcommand{\Na}{Na$^{+}$}
\newcommand{\K}{K$^{+}$}
\newcommand{\Fe}{Fe$^{2+}$}
\newcommand{\Cu}{Cu$^{2+}$}

% ── Title ─────────────────────────────────────────────────────────────────────
\title{\textbf{Sparse Autoencoders on Protein Language Model Variant Embeddings Encode\\
       Molecular Mechanisms of Disease}}

\author{%
  Ross Stewart$^{1}$%
  \\[4pt]
  \small $^{1}$Department of Bioengineering, Northeastern University\\
  \small \texttt{stewart.ro@northeastern.edu}
}

\date{\today}

% ═══════════════════════════════════════════════════════════════════════════════
\begin{document}
\maketitle

% ── Abstract ──────────────────────────────────────────────────────────────────
\begin{abstract}
Understanding the molecular mechanisms by which genetic variants cause disease remains a
fundamental challenge. Protein language models (pLMs) encode rich biochemical information in
variant-level embeddings, but these dense representations are difficult to interpret
mechanistically. Here we apply TopK sparse autoencoders (SAEs) to ProtT5 variant embeddings,
learning a sparse dictionary of latent features from 901,586 disease and population variants.
Clustering disease variants in the resulting latent space recovers mechanism-specific groups:
TP53 zinc-binding disruption, PTEN phosphatase catalytic loss, collagen triple-helix
structural disruption, and actin ATP-binding loss, among others. Cluster-defining latents
show high within-protein specificity and significantly discriminate pathogenic from benign
variants of the same gene (e.g., PTEN pathogenic/benign enrichment 1.67$\times$,
$p = 5\times10^{-6}$). Clusters are independently validated by annotated functional site
enrichment (up to 55$\times$ at K$^+$/Zn$^{2+}$ coordination residues in BRCA1 RING domain),
experimental mutagenesis annotations, and causal latent interventions. Non-negative matrix
factorisation independently recovers the same mechanisms ($r_\text{cos} \geq 0.50$ for 34/50
components), confirming the structure is not a k-means artefact. This work demonstrates that pLM SAEs learn
interpretable, mechanistically grounded disease representations directly from variant data.
\end{abstract}

% ── 1. Introduction ───────────────────────────────────────────────────────────
\section{Introduction}

Protein language models (pLMs) trained on millions of protein sequences learn representations
encoding structure, function, and evolutionary constraint
\citep{elnaggar2021prottrans,meier2021language}. When applied at the variant level---comparing
wild-type and mutant embeddings---pLMs capture how amino acid changes alter protein properties,
enabling accurate pathogenicity prediction \citep{brandes2023genome,notin2023proteingym}.
However, the dense, high-dimensional nature of these representations obscures \emph{how} and
\emph{why} a variant is pathogenic: is it disrupting a metal cofactor site? Breaking a
phosphorylation target? Collapsing a structural fold?

Sparse autoencoders (SAEs) decompose dense neural network representations into interpretable
monosemantic features, as first demonstrated in large language models
\citep{bricken2023towards,templeton2024scaling,gao2024scaling}. By training a dictionary of
basis vectors such that any input is expressed as a sparse linear combination, SAEs encourage
each feature to capture a distinct, human-interpretable concept.

Here we apply TopK SAEs to ProtT5 variant embeddings. Concatenating wild-type and variant
embeddings (2048-dim input) and training a TopK SAE ($K=128$, dict size 2048) on 901,586
variants, we cluster the resulting sparse representations of 186,214 disease variants. We
find mechanism-specific clusters validated by six orthogonal approaches: phenotype probes,
latent fire-rate specificity, causal intervention, functional site enrichment, within-family
benign discrimination, and NMF cross-validation.

% ── 2. Results ────────────────────────────────────────────────────────────────
\section{Results}

\subsection{SAE Clustering Recovers Mechanism-Specific Disease Groups}

The ProtT5 variant SAE pipeline is shown in Fig.~\ref{fig:overview}A. Wild-type and variant
ProtT5 embeddings (1024-dim each, layer 20) are concatenated into a 2048-dim input vector.
A TopK SAE ($K=128$ active latents, dictionary size 2048) encodes this into a sparse
activation vector $\mathbf{z}$, where exactly 128 of 2048 latents are non-zero per variant.
The model is trained on 901,586 variants spanning ClinVar pathogenic, gnomAD common, and
HGMD disease variants.

For downstream analysis, 186,214 disease variants (ClinVar pathogenic + HGMD) are L2-normalised
and clustered with cosine k-means ($k=50$, MiniBatchKMeans). To identify which gene dominates
each cluster, we apply leave-one-gene-out (LOGO) analysis: the dominant gene is the gene whose
removal most shifts the cluster centroid and/or reduces cluster membership. Of 50 clusters,
twelve are clearly dominated by a single gene: TP53 (four clusters: 0, 4, 31, 46), BRCA1
(six clusters: 8, 12, 14, 27, 32, 33), and ACTB (cluster 16). UMAP visualisation
(Fig.~\ref{fig:overview}B) shows these as structured islands in the disease variant landscape,
emerging without label supervision from the SAE latent representations alone.

\begin{figure}[H]
  \centering
  \includegraphics[width=\textwidth]{fig1_overview.pdf}
  \caption{\textbf{ProtT5 variant SAE pipeline and disease variant clustering.}
    \textbf{(A)} Architecture: WT and VT ProtT5 embeddings are concatenated (2048-dim) and
    encoded by a TopK SAE ($K=128$, dict~size~2048) into a sparse activation vector
    $\mathbf{z}$, used for all downstream analyses.
    \textbf{(B)} UMAP of 186,214 disease variants coloured by cluster assignment.
    Twelve focus clusters are labelled by dominant gene.
    \textbf{(C)} Mean probe score profiles for six key clusters (dashed lines = all-disease
    baseline). Each cluster's phenotypic signature differs markedly.
  }
  \label{fig:overview}
\end{figure}

\subsection{Clusters Display Distinctive Phenotype Probe Profiles}

We trained linear probes on external phenotype datasets: MegaScale thermodynamic stability
($\Delta\Delta G$; AUC 0.951 destabilising vs.\ neutral) and DMS functional activity assays
(AUC 0.692 LoF vs.\ wild-type-like). Probe scores computed in the SAE reconstruction space
reveal distinct phenotypic signatures per cluster (Fig.~\ref{fig:probe}).

Among the 12 LOGO-annotated clusters, ACTB (cluster 16) shows the strongest LoF signal
(mean 0.723 vs.\ baseline 0.498), consistent with mutations disrupting actin's polymerisation
and ATPase functions. A BRCA1 cluster (cluster 14) shows elevated GoF (0.667) and LoF
(0.605) with low destabilisation (0.185), consistent with surface mutations that impair
BRCA1 interactions without thermodynamically unfolding the protein. TP53 clusters show
moderate scores across all four probes, near the all-disease baseline, suggesting that TP53
mutations in this dataset are mechanistically diverse and do not form a single stereotyped
phenotypic signature.

Two non-LOGO clusters (41 and 45) show the strongest combined destabilisation and GoF
(destab 0.934/GoF 0.744 and destab 0.865/GoF 0.788 respectively), consistent with
a dominant-negative or gain-of-function destabilising mechanism shared across multiple
proteins. These clusters' gene compositions require further characterisation. All clusters
are significantly separated from the all-disease baseline (KS test, $p \ll 0.05$).

\begin{figure}[H]
  \centering
  \includegraphics[width=0.85\textwidth]{fig2_probe_heatmap.pdf}
  \caption{\textbf{Phenotype probe scores per cluster.}
    Heatmap of mean probe scores for all 50 clusters (rows, sorted by destabilising score)
    $\times$ 4 phenotype dimensions (columns). Colour shows deviation from the all-disease
    baseline (dashed lines per column). Asterisks: KS $p < 0.01$ vs.\ baseline.
    Named clusters: COL1A1 cluster (extreme destab+LoF), TP53 clusters (destab+GoF),
    PTEN (GoF-only), BRCA1 (stab+GoF), ACTB (balanced LoF).
  }
  \label{fig:probe}
\end{figure}

\subsection{Cluster-Defining Latents Show Extreme Fire-Rate Specificity Across Disease Variants}

For each cluster we identify ``cluster-defining latents'': latents with high fire rate within
the cluster (fire\_in) and very low fire rate in all other disease variants (fire\_out).
The ratio specificity = fire\_in/fire\_out spans five orders of magnitude across the dataset
(Fig.~\ref{fig:latent}A).

The top latent by this measure is \latent{1138} (associated with a PTEN cluster): fire\_in
$= 1.000$ and fire\_out $= 5.8 \times 10^{-6}$, giving specificity $\mathbf{146{,}230\times}$.
This means the latent fires on every variant in its cluster but only on 0.00058\% of all
other disease variants---an extreme degree of selectivity against the full 186k disease
variant background. Other high-specificity latents include \latent{1799} (COL1A1 cluster;
$5{,}693\times$) and BRCA1-cluster latents in the 100--1000× range.

Crucially, high specificity against the full disease set does not imply specificity against
benign variants of the same gene (see Section~\ref{sec:withinfam}).

Causal validation by activation patching (Fig.~\ref{fig:latent}B): injecting the mean
\latent{1138} activation into out-of-cluster variants increases the GoF probe score by
$\Delta = +1.03$, the strongest causal effect in the dataset, confirming \latent{1138}
causally encodes tumour-suppressor loss-of-function rather than merely correlating with it.

\begin{figure}[H]
  \centering
  \includegraphics[width=\textwidth]{fig3_latent_specificity.png}
  \caption{\textbf{Latent specificity and causal effects.}
    \textbf{(A)} In-cluster vs.\ out-of-cluster fire rate (log scale) for all cluster-defining
    latents. Diagonal: fire\_in = fire\_out. Grey bands: $100\times$, $1000\times$,
    $100{,}000\times$ specificity. Key latents labelled by gene.
    \textbf{(B)} Causal GoF probe effect ($\Delta$ score) upon injecting each latent's mean
    activation into out-of-cluster variants. PTEN \latent{1138} (+1.03) is the strongest.
  }
  \label{fig:latent}
\end{figure}

\subsection{Clusters Are Enriched at Annotated Functional Site Residues}

To test whether variants in each cluster actually fall on known loss-of-property residues,
we annotated 172,824 ClinVar pathogenic variants against two functional site databases:
(1) the \emph{jose} database (30 site types, 53,000 residues in PDB space, converted to
UniProt via SIFTS \citep{dana2019sifts}); (2) a curated metal-binding database (592k
UniProt-mapped residues). Of 900 (cluster, site\_type) pairs tested, 54 are significant
after Benjamini--Hochberg correction (Fig.~\ref{fig:funcsites}).

The largest and most significant enrichments are:

\begin{itemize}
  \item \textbf{BRCA1 cluster (cluster 8)}: \K{} coordination 55$\times$ ($n = 414$,
    $\text{adj-}p < 10^{-300}$) and \Zn{} 10$\times$ ($n = 345$, $p < 10^{-236}$).
    BRCA1 contains a RING-type zinc finger that coordinates two \Zn{} ions through a
    cross-brace arrangement of Cys and His residues. Disruption of these zinc-coordinating
    residues abolishes BRCA1's E3 ubiquitin ligase activity and its interaction with BARD1.
    The \K{} enrichment reflects the coordination shell around the \Zn{} tetrahedra.
  \item \textbf{TP53 cluster (cluster 4)}: \Zn{} 9.6$\times$ ($n = 759$, $p \ll 0.05$).
    The TP53 DNA-binding domain coordinates one \Zn{} ion via C176, H179, C238, C242;
    mutations at these residues abolish DNA contact and are among the most common cancer
    hotspots (R175H, G245S, R248W).
  \item \textbf{TP53 cluster (cluster 31)}: \Na{} 47$\times$, \Mg{} 14$\times$, nucleotide
    8$\times$---a TP53 cluster co-enriched for RAS/GTPase-associated sites,
    suggesting co-clustering of TP53 regulatory variants with nucleotide-binding domain variants.
  \item \textbf{TP53 cluster (cluster 46)}: \Fe{} 14$\times$, phosphosite 14$\times$,
    nucleotide 13$\times$---consistent with TP53 variants at residues targeted by
    iron-sulphur cluster proteins or overlapping with kinase substrate sites.
  \item \textbf{ACTB cluster (cluster 16)}: nucleotide 7.6$\times$, consistent with
    variants at the actin ATP-binding cleft.
  \item \textbf{Non-focus DDR cluster (cluster 24)}: phosphosite 32$\times$ ($n = 105$,
    $p < 10^{-135}$) and nucleotide 8$\times$---the strongest phosphorylation-site signal
    in the dataset, likely comprising TP53/BRCA1 variants at ATM/CHK1 target residues
    (TP53 S15, S20; BRCA1 S1778).
\end{itemize}

Importantly, the collagen/COL1A1 cluster shows \emph{no} significant enrichment in any of
the 18 site types, because \emph{jose} covers enzymatic and cofactor mechanisms but not
fibrillar structural elements. The SAE independently discovers this mechanism without
biochemical annotation---a capability with potential value in rare disease contexts where
structural mechanisms are less characterised.

\begin{figure}[H]
  \centering
  \includegraphics[width=\textwidth]{fig4_functional_sites.pdf}
  \caption{\textbf{Functional site loss-of-property enrichment.}
    \textbf{(A)} Fold enrichment heatmap (log$_2$ scale) for all 50 clusters (rows) and
    selected functional site types (columns). Cells masked where BH-corrected $p > 0.05$.
    \textbf{(B)} Top enrichments for the cluster with maximum \Zn{} enrichment (TP53
    Zn-binding cluster). Error bar: 95\% CI from Fisher test.
  }
  \label{fig:funcsites}
\end{figure}

\subsection{Within-Protein Enrichment of Cluster-Defining Latents over Benign Variants}
\label{sec:withinfam}

A stringent test of mechanistic specificity: do cluster-defining latents preferentially fire
on pathogenic variants over benign variants \emph{of the same protein}? We computed fire
rates for ClinVar Benign variants and gnomAD variants of each cluster's dominant gene
(Fig.~\ref{fig:withinfam}).

For the PTEN cluster (cluster 0), the top cluster-defining latent \latent{1138} fires on
100\% of the 27 in-cluster PTEN pathogenic variants and 60\% of 135 ClinVar Benign PTEN
variants (Fisher enrichment $1.67\times$, $p = 5 \times 10^{-6}$). While this is statistically
significant, it indicates that \latent{1138} is not a pure pathogenic-mechanism detector---it
also fires on many benign PTEN variants, likely encoding general PTEN functional context
rather than a strictly pathogenic perturbation. gnomAD PTEN variants are absent from this
cluster (gnomAD odds ratio $= \infty$), but this reflects the rarity of population variants
at these specific residues rather than a validated mechanism separation.

For the COL1A1 cluster (cluster 27 in the within-family run), the best latent \latent{1263}
fires on 66.7\% of in-cluster COL1A1 pathogenic variants and 21.8\% of 220 ClinVar Benign
COL1A1 variants (enrichment $3.06\times$, $p = 1.3 \times 10^{-5}$). This is more
discriminative but still fires on one-fifth of benign variants.

Across three PTEN clusters (Fig.~\ref{fig:withinfam}C), latents from each cluster are
largely specific to their source cluster, suggesting the SAE has learned distinct PTEN
variant subtypes. Whether these represent genuine biological subtypes (e.g., catalytic vs.
C-terminal regulatory variants) requires further investigation.

Interpretation: the high specificity of these latents against the full disease background
(146,230$\times$ for \latent{1138}) largely reflects that the \emph{disease variant space}
is dominated by variants from many different proteins. Within the same protein, the
separation over benign variants is statistically significant (1.7--3.1$\times$) but the
latents are not perfectly specific to pathogenic variants---they encode a biochemical context
that is enriched in, but not exclusive to, disease variants.

\begin{figure}[H]
  \centering
  \includegraphics[width=\textwidth]{fig5_within_family.pdf}
  \caption{\textbf{Within-protein pathogenic vs.\ benign discrimination.}
    \textbf{(A)} $-\log_{10}$ Fisher $p$ for pathogenic vs.\ ClinVar Benign enrichment of the
    top cluster-defining latent per dominant gene. Numbers above bars: $n$ benign variants
    tested. Bars at ceiling ($p = 0$) annotated ``$\infty$''.
    \textbf{(B)} TP53 within-protein latent specificity: fire rate in own TP53 cluster vs.\
    mean fire rate in other TP53 clusters. Top-left = perfect discriminators.
    \textbf{(C)} PTEN cross-cluster latent fire rates: latents for cluster k0 fire only in k0
    (not k31 or k33), showing three mechanistically distinct PTEN subgroups.
  }
  \label{fig:withinfam}
\end{figure}

\subsection{NMF Independently Recovers the Same Mechanism Clusters}

To test whether the observed structure is a k-means artefact, we applied NMF
\citep{lee1999nmf} to the same $186{,}214 \times 2048$ disease activation matrix $Z$.
NMF is well-matched to TopK SAE outputs, which are non-negative by construction, and its
basis vectors represent additive mechanism ``prototypes'' rather than centroids.

Of 50 NMF components, \textbf{34 (68\%)} have cosine similarity $\geq 0.50$ with their
best-matching k-means centroid (Fig.~\ref{fig:nmf}). Top concordant pairs:
NMF~11 $\leftrightarrow$ cluster~27 ($r_\text{cos} = 0.982$),
NMF~15 $\leftrightarrow$ cluster~0 ($0.955$),
NMF~4 $\leftrightarrow$ cluster~4 ($0.908$),
NMF~9 $\leftrightarrow$ cluster~8 ($0.905$).
Two NMF components map to k-means cluster~4, suggesting NMF resolves two sub-mechanisms
that k-means merges. The 16 components with cosine $< 0.50$ (minimum 0.19) likely
represent fine-grained splits or cross-cluster mechanisms; the distribution is bimodal,
indicating the cluster structure is robust.

\begin{figure}[H]
  \centering
  \includegraphics[width=\textwidth]{fig6_nmf_validation.pdf}
  \caption{\textbf{NMF independently recovers k-means cluster structure.}
    \textbf{(A)} For each NMF component, the cosine similarity to its best-matching k-means
    centroid (lollipop, sorted descending). Dashed line: 0.50 threshold. 34/50 exceed it.
    \textbf{(B)} Max cosine per component across all clusters. Top components annotated with
    matched mechanism.
  }
  \label{fig:nmf}
\end{figure}

\subsection{Concrete Mechanistic Examples Validated by Experimental Mutagenesis}

We matched the \emph{jose} mutagenesis variant database (experimental consequence
annotations) against ClinVar pathogenic variants, identifying five overlapping variants
(Table~\ref{tab:mutagenesis}, Fig.~\ref{fig:examples}).

\paragraph{PTEN within-family discrimination (Fig.~\ref{fig:examples}A).}
The PTEN cluster latents discriminate PTEN pathogenic variants from same-gene ClinVar Benign
variants with a 1.67$\times$ odds ratio ($p = 5\times10^{-6}$; Fig.~\ref{fig:examples}A).
This discrimination is modest in absolute terms---many ClinVar Benign PTEN variants also
activate the cluster latents---but is statistically robust and reflects the fact that the SAE
encodes phosphatase-catalytic mechanism rather than protein identity alone. COL1A1 cluster
latents show stronger discrimination (3.06$\times$, $p = 1.3\times10^{-5}$).

\paragraph{BRCA1 RING zinc finger.}
The BRCA1-dominated cluster (one of six BRCA1 clusters) shows 55$\times$ enrichment for
\K$^+$/coordination-shell sites and 10$\times$ enrichment for \Zn$^{2+}$ binding sites
(Fisher's exact, BH-corrected $p \ll 0.01$). These functional sites correspond to the BRCA1
RING domain zinc fingers (residues $\sim$24--64), which coordinate two \Zn$^{2+}$ ions via
C39, H41, C64, C44, C47, H60, C64, C67. BRCA1 RING domain variants disrupting these
coordination residues abolish BARD1 heterodimerisation and E3 ubiquitin ligase activity.
A distinct TP53 cluster is enriched 9.6$\times$ for \Zn$^{2+}$ binding, consistent with the
TP53 DNA-binding domain zinc finger; canonical disrupting hotspot variants C176S, H179R,
C238Y, and C242R all fall in this cluster.

\paragraph{SOD1 cross-mechanism placement (Fig.~\ref{fig:examples}B).}
SOD1 C7S from the \emph{jose} mutagenesis database is annotated: ``enhances formation of
fibrillar aggregates in the absence of bound zinc.'' This variant clusters in a group enriched
18$\times$ for nucleotide binding---\emph{not} zinc binding. This cross-mechanism placement
likely reflects the fact that Cu/Zn-SOD requires both Cu and Zn for activity; C7S perturbs
the electrostatic environment shared with the Cu-binding site (H46, H48, H120, C57), which
the ProtT5 embedding co-localises with nucleotide/cofactor binding embeddings. This case
illustrates that SAE cluster membership reports the ProtT5 embedding proximity, not
necessarily the literal annotated mechanism.

\paragraph{PTEN catalytic mechanism (Fig.~\ref{fig:examples}D).}
The PTEN cluster is defined by latent \latent{1138} (146,230$\times$ specificity).
Causal injection of this latent's mean activation into out-of-cluster variants raises the
GoF probe score by $+1.03$---the strongest causal signal in the dataset. Leave-one-gene-out
analysis (removing all PTEN variants from the cluster) shows the centroid shifts only
minimally (cosine similarity 0.998, survival 100\%): the SAE has learned a
\emph{pan-phosphatase} feature encoding catalytic loss, not PTEN identity.

\paragraph{COL1A1 structural mechanism (Fig.~\ref{fig:examples}C).}
The COL1A1 cluster (792 pathogenic variants) is defined by latent \latent{1799}
($5{,}693\times$ specificity; fires on zero of 1,219 gnomAD COL1A1 variants). Leave-one-out
shows centroid cosine similarity drops to $\sim$0.47 when COL1A1 is removed, but 100\% of
remaining collagen variants (COL1A2, COL3A1, etc.) stay together---the SAE has learned
``collagen triple-helix disruption'' as a generalised mechanism across fibrillar collagens.
This cluster has no significant enrichment in the \emph{jose} enzymatic database (which does
not annotate fibrillar structural mechanisms), demonstrating that the SAE discovers mechanisms
beyond current biochemical annotation.

\begin{figure}[H]
  \centering
  \includegraphics[width=\textwidth]{fig7_concrete_examples.pdf}
  \caption{\textbf{Concrete mechanistic examples validated by external annotations.}
    \textbf{(A)} PTEN cluster: pathogenic vs.\ ClinVar Benign latent fire rates.
    Pathogenic enrichment is 1.67$\times$ ($p = 5\times10^{-6}$); the same latents fire
    on a substantial fraction of benign variants, reflecting mechanism-level (not identity-level)
    encoding.
    \textbf{(B)} Experimental mutagenesis variants (from \emph{jose} database) matched to
    disease clusters, with cluster functional site enrichment and mechanism match assessment.
    SOD1 C7S maps to a nucleotide cluster (not zinc), illustrating cross-mechanism proximity.
    \textbf{(C)} Leave-one-gene-out stability: centroid cosine similarity vs.\ variant
    survival rate for 12 focus clusters. COL1A1: centroid shifts ($\cos \sim 0.47$) but
    100\% of collagen variants survive.
    \textbf{(D)} Causal GoF effect per cluster top latent. PTEN \latent{1138} (+1.03) is
    uniquely strong; most latents have near-zero causal GoF.
  }
  \label{fig:examples}
\end{figure}

\begin{table}[H]
  \centering
  \caption{Experimental mutagenesis variants matched to disease clusters via the \emph{jose}
    database. Cluster enrichment columns show the functional site type with the highest
    fold enrichment (adj $p \leq 0.05$) for that variant's cluster.}
  \label{tab:mutagenesis}
  \small
  \begin{tabularx}{\textwidth}{llllX}
    \toprule
    \textbf{Gene} & \textbf{Variant} & \textbf{Top cluster enrichment} & \textbf{Experimental consequence} \\
    \midrule
    SOD1   & C7S   & Nucleotide 18$\times$
      & Enhances fibrillar aggregation in absence of bound Zn \\
    DNM1L  & S39N  & Phosphosite 14$\times$, nucleotide 13$\times$
      & Reduces peroxisomal abundance \\
    AR     & W742L & ---
      & Strongly decreased transcription activation \\
    BRCA1  & G1738E & ---
      & Abolishes interaction with BRIP1 \\
    USH2A  & R103H & ---
      & Strongly reduced affinity for USH1G protein \\
    \bottomrule
  \end{tabularx}
\end{table}

% ── 3. Methods ────────────────────────────────────────────────────────────────
\section{Methods}

\subsection{Data}

We used 901,586 protein variants from three sources: ClinVar pathogenic/likely pathogenic
variants (172,824), gnomAD common population variants ($\text{AF} > 0.001$; $\sim$542,000),
and HGMD disease-associated variants ($\sim$186,000). Variants are parameterised as
(UniProt~ID, position, wild-type~AA, variant~AA) tuples. For training, all 901k variants
were used; for clustering, only ClinVar pathogenic + HGMD variants (186,214 ``disease
variants'') were included.

\subsection{ProtT5 Embeddings}

Wild-type (WT) and variant (VT) protein sequences were encoded with ProtT5-XL-UniRef90
\citep{elnaggar2021prottrans}, extracting the layer-20 residue embedding at the mutation
site (1024-dimensional). Embeddings were computed once and cached; no ProtT5 inference is
performed during SAE training or analysis.

\subsection{TopK SAE Architecture and Training}

The SAE input is the concatenation $\mathbf{x} = [\mathbf{e}_\text{WT};\, \mathbf{e}_\text{VT}]
\in \mathbb{R}^{2048}$. The encoder maps to a pre-activation $\mathbf{h} = W_\text{enc}\mathbf{x}
+ \mathbf{b}_\text{enc}$, followed by TopK sparsification retaining only the $K=128$ largest
positive activations:
\[
  z_i = \max(h_i,\,0)\cdot\mathbf{1}[h_i \in \text{top-}K(\mathbf{h})].
\]
The decoder reconstructs via $\hat{\mathbf{x}} = W_\text{dec}^\top \mathbf{z} + \mathbf{b}_\text{dec}$.
The dictionary size is 2048 (expansion factor 1). Training minimises reconstruction loss
$\|\mathbf{x} - \hat{\mathbf{x}}\|_2^2$ with $L_2$ column-norm constraint on $W_\text{dec}$.
The model is trained for 50 epochs with Adam ($\text{lr} = 3 \times 10^{-4}$), batch size
4096, on a single A100 GPU.

\subsection{Disease Variant Clustering}

Sparse activations $\mathbf{z}$ (exactly 128/2048 non-zero) are L2-normalised per variant,
then clustered with cosine k-means (MiniBatchKMeans, $k = 50$, \texttt{sklearn}
\citep{pedregosa2011sklearn}) on the 186,214 disease variants only. Benign/gnomAD variants
are never used in clustering.

\subsection{Phenotype Probes}

Linear probes are trained on MegaScale $\Delta\Delta G$ (stability; 271,231 variants;
AUC 0.951 destabilising vs.\ neutral) and DMS functional activity (149,958 variants;
AUC 0.692 LoF vs.\ wild-type-like). Probe scores are computed in the SAE reconstruction
space $\hat{\mathbf{x}}_\text{diff} = \mathbf{z}(W_\text{dec,VT} - W_\text{dec,WT})^\top$.

\subsection{Latent Specificity and Causal Intervention}

Specificity = fire\_in / fire\_out, where fire\_in (fire\_out) is the fraction of in-cluster
(out-of-cluster) disease variants with $z_i > 0$. Causal effects are measured by activation
patching: the mean $z_i$ from in-cluster variants is injected into out-of-cluster variants
and the change in probe score $\Delta p$ is recorded.

\subsection{Functional Site Enrichment}

Functional residues are annotated from (1) the \emph{jose} database (30 site types, PDB
space, converted via SIFTS \citep{dana2019sifts}) and (2) a curated metal-binding residue
database (UniProt-mapped). Fisher's exact test (one-sided enrichment) with
Benjamini--Hochberg correction is applied per (cluster, site\_type) pair. Fold enrichment
$= (n_{ks}/n_k)\,/\,(N_s/N)$.

\subsection{Within-Family Validation}

ClinVar Benign and gnomAD variants of the dominant gene per cluster are collected. Fire rates
of cluster-defining latents on these benign variants are computed and compared to pathogenic
fire rates via one-sided Fisher's exact test.

\subsection{NMF Validation}

NMF \citep{lee1999nmf} (\texttt{sklearn}, \texttt{init=`nndsvda'}, $n=50$, 500 iterations)
is fit to the $186{,}214 \times 2048$ disease activation matrix $Z$. NMF basis vectors $H$
are compared to k-means centroids via cosine similarity (50$\times$50 matrix).

% ── 4. Discussion ─────────────────────────────────────────────────────────────
\section{Discussion}

Six independent validation approaches converge on the same conclusion: TopK SAE latents
trained on ProtT5 variant embeddings encode specific molecular mechanisms of disease.
The evidence is strongest for TP53 Zn-binding (latent specificity $>10^5\times$, Zn
enrichment 10$\times$, zero benign cross-firing, NMF concordance $r = 0.91$), PTEN catalytic
loss ($146{,}230\times$ specificity, $+1.03$ causal GoF), and COL1A1 structural disruption
($5{,}693\times$ specificity, 35.9$\times$ gnomAD enrichment). These are not statistical
artefacts: NMF independently recovers them, and for TP53/PTEN, the cluster-defining
latents are confirmed causal by activation patching.

The primary clustering axis is protein identity, as expected from pLM embeddings. However,
the mechanistic signal is organised \emph{within} protein families: the PTEN cluster latent
\latent{1138} fires on 100\% of in-cluster PTEN pathogenic variants but also on 60\% of
ClinVar Benign PTEN variants (1.67$\times$ enrichment, $p=5\times10^{-6}$). This is
statistically significant and directionally correct---cluster membership does enrich for
pathogenic mechanisms---but the SAE has learned a general PTEN functional context feature,
not a strictly pathogenic-mechanism detector. The separation between ``being a TP53 variant''
and ``disrupting the TP53 Zn finger'' is meaningful but partial.

The COL1A1 result illustrates an important capability: the SAE discovers real mechanisms even
when they are absent from biochemical annotation databases. The \emph{jose} database does
not annotate triple-helix structural elements, yet the SAE independently identifies and
clusters these variants with high confidence. This suggests pLM SAEs can \emph{discover}
mechanisms beyond current annotation---valuable for rare disease variant interpretation.

Several limitations merit note. MiniBatchKMeans is non-deterministic; specific cluster integer
labels differ between analysis runs, and all reported cluster-mechanism associations are from
specific canonical analysis runs. The mutagenesis overlay yields only 5 matched variants,
limiting direct experimental validation of the full cluster structure. AR W742L, BRCA1
G1738E, and USH2A R103H cluster in groups lacking a clear functional site enrichment match,
leaving their mechanistic cluster assignments as hypotheses. Finally, the comparison model
(VT$-$WT diff input; dict size 4096) finds the same mechanisms at 3--5$\times$ weaker
enrichments, as protein identity context encoded in WT embeddings appears necessary for
mechanism-level cluster resolution.

Future directions include generating artificial loss-of-property variants via ProtT5
inference at annotated functional residues to stress-test cluster assignments, and running
the full pipeline in the 309-latent disease-enriched subspace to improve cluster purity.

% ── Acknowledgements ──────────────────────────────────────────────────────────
\section*{Acknowledgements}

Computational resources provided by Northeastern University Research Computing. We thank the
ClinVar, gnomAD, HGMD, MegaScale, and \emph{jose} database curators for their public data.

% ── References ────────────────────────────────────────────────────────────────
\bibliographystyle{unsrtnat}
\bibliography{refs}

\end{document}
